%% kapa: exam=1 round=34 year=2023 period=2 form=A subject=accounting subject_ko=회계학 seq=20 number=60
%% source: Q-Net 감정평가사 기출문제 (cst00309) artlSeq=5211789 file=raw/1008/5211789/2023년도 제34회 감정평가사 1차시험 2교시.pdf
%% conversion: kordoc 4.24.51 -> tools/gen1.py (auto); review=proofread-sonnet; flags=none
\begin{question}{60}
(주)감평은 기계장치(장부금액 ￦2,000, 공정가치 ￦3,500)를 제공하고, (주)한국의 건물과 현금 ￦700을 취득하는 교환거래를 하였다. 건물의 공정가치는 ￦2,500으로 기계장치의 공정가치보다 더 명백하며, 이 교환거래는 상업적 실질이 있다고 할 때, (주)감평이 인식할 유형자산처분손익은?
\begin{choices}
  \item 유형자산처분손익 ￦0
  \item 유형자산처분손실 ￦1,200
  \item 유형자산처분이익 ￦1,200
  \item 유형자산처분손실 ￦2,200
  \item 유형자산처분이익 ￦2,200
\end{choices}
\end{question}
\answer{3}
